\chapter{CONCLUSIONES}

\begin{enumerate}
    \item Se formuló rigurosamente la métrica y geodésicas de Kerr-Newman con carga débil en unidades SI y firma $(+,-,-,-)$, confirmando que la carga disminuye el radio crítico de la esfera de fotones.
    \item Se derivó paso a paso la ecuación del oscilador inestable para el exponente de Lyapunov $\gamma_\lambda$, justificando la divergencia exponencial de las trayectorias vecinas.
    \item Se demostró que el coeficiente de segundo orden de Kobialko $D = 1/27$ constituye la huella analítica única para distinguir la carga de Kerr-Newman frente a soluciones regulares de Bardeen y Hayward.
    \item Se comprobó que el ensanchamiento de Sneppen en métricas rotantes incrementa la visibilidad interferométrica de los sub-anillos $n \ge 1$ para telescopios de nueva generación.
\end{enumerate}
